% Trecho de demonstração. Texto conferido à mão contra a imagem da página.
% Marcadores %%>>> id / %%<<< id delimitam trechos reaproveitáveis nos slides.
\section*{O ente immutavel}

%%>>> ente-immutavel
\pg{37}\textbf{87. O ente immutavel é necessario, infinito, simplicissimo.} ---
O ente immutavel é acto puro. Ora o ente, que é acto puro, é necessario, infinito,
simplicissimo. Logo o ente immutavel é necessario, infinito, simplicissimo.

O ente necessario, infinito, simplicissimo, é só acto, sem nenhuma sombra de
potencialidade.\pg{38} Ora o que é só acto, sem sombra de potencialidade, é
immutavel. Logo o ente necessario, infinito, simplicissimo, é immutavel.
%%<<< ente-immutavel

\begin{latim}
Primum ens oportet esse purum actum absque permixtione alicujus potentiae, eo quod
potentia simpliciter est posterior actu. \textup{(Sum. Th., I, q.\,9, a.\,1)}
\end{latim}
